\documentclass[10pt,a4paper]{amsart}
\usepackage{amsmath,amssymb,amsthm,graphicx}
% Retained minimal preamble for the Stacks Project chapter more-algebra.tex (More on Algebra).
% The chapter file pinned-original-source.tex carries only \input{preamble} on its first line, so no
% preamble text lives inside the chapter itself.  The Stacks Project uses one single shared preamble for
% every chapter; the snapshot stages it identically as sources/stacks-chow/preamble.tex and
% sources/stacks-etale-cohomology/preamble.tex.  The block below is copied mechanically, line for line,
% from sources/stacks-chow/preamble.tex lines 14-18 (the Xy-pic setup) and 154-217 (the native theorem
% declarations together with the mathematical macros).  Nothing was added, removed or redefined; the
% wrapper class and the amsmath/amssymb/amsthm/graphicx base packages are supplied by the assembler.
% Chapter: more-algebra.tex -- More on Algebra.
% For commutative diagrams we use Xy-pic
\usepackage[all]{xy}

% We use 2cell for 2-commutative diagrams.
\xyoption{2cell}

% Theorem environments.
%
\theoremstyle{plain}
\newtheorem{theorem}[subsection]{Theorem}
\newtheorem{proposition}[subsection]{Proposition}
\newtheorem{lemma}[subsection]{Lemma}

\theoremstyle{definition}
\newtheorem{definition}[subsection]{Definition}
\newtheorem{example}[subsection]{Example}
\newtheorem{exercise}[subsection]{Exercise}
\newtheorem{situation}[subsection]{Situation}

\theoremstyle{remark}
\newtheorem{remark}[subsection]{Remark}
\newtheorem{remarks}[subsection]{Remarks}

\numberwithin{equation}{subsection}

% Macros
%
\def\lim{\mathop{\mathrm{lim}}\nolimits}
\def\colim{\mathop{\mathrm{colim}}\nolimits}
\def\Spec{\mathop{\mathrm{Spec}}}
\def\Hom{\mathop{\mathrm{Hom}}\nolimits}
\def\Ext{\mathop{\mathrm{Ext}}\nolimits}
\def\SheafHom{\mathop{\mathcal{H}\!\mathit{om}}\nolimits}
\def\SheafExt{\mathop{\mathcal{E}\!\mathit{xt}}\nolimits}
\def\Sch{\mathit{Sch}}
\def\Mor{\mathop{\mathrm{Mor}}\nolimits}
\def\Ob{\mathop{\mathrm{Ob}}\nolimits}
\def\Sh{\mathop{\mathit{Sh}}\nolimits}
\def\NL{\mathop{N\!L}\nolimits}
\def\CH{\mathop{\mathrm{CH}}\nolimits}
\def\proetale{{pro\text{-}\acute{e}tale}}
\def\etale{{\acute{e}tale}}
\def\QCoh{\mathit{QCoh}}
\def\Ker{\mathop{\mathrm{Ker}}}
\def\Im{\mathop{\mathrm{Im}}}
\def\Coker{\mathop{\mathrm{Coker}}}
\def\Coim{\mathop{\mathrm{Coim}}}

% Boxtimes
%
\DeclareMathSymbol{\boxtimes}{\mathbin}{AMSa}{"02}

%
% Macros for moduli stacks/spaces
%
\def\QCohstack{\mathcal{QC}\!\mathit{oh}}
\def\Cohstack{\mathcal{C}\!\mathit{oh}}
\def\Spacesstack{\mathcal{S}\!\mathit{paces}}
\def\Quotfunctor{\mathrm{Quot}}
\def\Hilbfunctor{\mathrm{Hilb}}
\def\Curvesstack{\mathcal{C}\!\mathit{urves}}
\def\Polarizedstack{\mathcal{P}\!\mathit{olarized}}
\def\Complexesstack{\mathcal{C}\!\mathit{omplexes}}
% \Pic is the operator that assigns to X its picard group, usage \Pic(X)
% \Picardstack_{X/B} denotes the Picard stack of X over B
% \Picardfunctor_{X/B} denotes the Picard functor of X over B
\def\Pic{\mathop{\mathrm{Pic}}\nolimits}
\def\Picardstack{\mathcal{P}\!\mathit{ic}}
\def\Picardfunctor{\mathrm{Pic}}
\def\Deformationcategory{\mathcal{D}\!\mathit{ef}}

% --- appended 2026-10-12 (Chief): the project own definitions of the reference/slogan/history
% display environments, copied mechanically from the shared Stacks preamble lines 8-12
% (\usepackage{verbatim} plus the three \newenvironment lines, which the chapter preamble keeps
% verbatim from the project preamble). Nothing added, removed or redefined: these environments
% carry no mathematical content and the project itself typesets them as comments. Without these
% three lines a counted statement whose own text uses them fails with Environment slogan undefined.
\usepackage{verbatim}
\newenvironment{reference}{\comment}{\endcomment}
\newenvironment{slogan}{\comment}{\endcomment}
\newenvironment{history}{\comment}{\endcomment}
\usepackage{cleveref}
\crefname{originallemma}{original lemma}{original lemmas}
\Crefname{originallemma}{Original lemma}{Original lemmas}
\makeatletter
\newlabel{lemma-ideal-well-defined}{{27893}{source}{}{sourcepointer-1}{}}
\newlabel{lemma-ideal-well-defined@cref}{{[originallemma][0][]27893}{[1][0][]source}{}{}{}}
\makeatother
\title{For principal Fitting-type ideals the beta-invariant is independent of the chosen representative complex}
\author{The Stacks Project Authors}
\date{stacks-project:more-algebra.tex:local-snapshot-20261012; GNU FDL 1.2 or later}
\begin{document}
\maketitle
\setcounter{subsection}{1}
\noindent\textit{Editorial note.} All delivered mathematical text is the author's own, copied from the pinned Stacks Project chapter file more-algebra.tex; no statement or proof step was rewritten, reordered or removed.\par\medskip
\noindent\textit{Editorial note.} The chapter file pinned-original-source.tex carries only \textbackslash{}input\{preamble\} on its first line, so no preamble text lives inside the chapter. The retained file sources/stacks-more-algebra/retained-preamble.tex is a mechanical, line-for-line copy of the project's single shared preamble as staged in the snapshot; nothing was added, removed or redefined and no package was added.\par\medskip
\noindent\textit{Editorial note.} The source URL is the source-level Stacks Project URL, as required when more than one card is produced from one source; the per-claim tag is carried by claim\_key.\par\medskip
\noindent\textit{Editorial note.} The author's first sentence reads 'Observe that the fact that I\_i(M\textasciicircum{}bullet, f) is a principal ideal implies that I\_i(M\textasciicircum{}bullet, f) is a principal ideal by Lemma ...', stating the same ideal twice; that is the author's text and is reproduced unchanged. The author closes with 'This is a computation that we omit.', also kept verbatim. The two references inside the counted proof span are recorded as non-load-bearing navigation pointers.\par\medskip
\noindent\textit{Allowed standard prerequisites (editorial).}\begin{itemize}
\item The ideals I\_i and J\_i attached to a bounded complex of finite free modules and a nonzerodivisor, and their universal property. used throughout the counted proof span. Fixes the two ideals whose equality is asserted and the property used to reduce to summands.
\item Well-definedness of I\_i up to powers of f under change of representative complex. invoked at proof lines 28175-28178. Makes the principality hypothesis meaningful and supplies the reduction to a direct summand.
\item Compatibility of the eta transform with direct sums. used at proof lines 28183-28189. Is the step that splits the displayed map into the contributions of M\textasciicircum{}bullet and Q\textasciicircum{}bullet.
\end{itemize}\medskip
\ifdefined\hypertarget\hypertarget{sourcepointer-1}{}\fi
\noindent\textit{Source reference.} Original lemma 27893: sources/stacks-more-algebra/pinned-original-source.tex, source lines 27893--27893. Same-chapter navigation pointer; the statement it names is standard for the section and the argument that consumes it stays inside the counted proof span.\par\medskip
\begin{lemma}
\label{lemma-eta-vanishing-beta-plus-pre}
Let $A$ be a ring and let $f \in A$ be a nonzerodivisor. Let
$M^\bullet$ and $N^\bullet$ be two bounded complexes of finite
free $A$-modules representing the same object in $D(A)$.
Assume $I_i(M^\bullet, f)$ is a principal ideal
for all $i \in \mathbf{Z}$. Then $J_i(M^\bullet, f) = J_i(N^\bullet, f)$
as ideals in $A/fA$.
\end{lemma}

\begin{proof}
Observe that the fact that $I_i(M^\bullet, f)$ is a principal ideal
implies that $I_i(M^\bullet, f)$ is a principal ideal by
Lemma \ref{lemma-ideal-well-defined} and hence the statement makes sense.
As in the proof of Lemma \ref{lemma-ideal-well-defined}
we may assume $N^\bullet = M^\bullet \oplus Q^\bullet$
for some trivial complex $Q^\bullet$, i.e.,
$$
Q^\bullet = \ldots \to 0 \to A \xrightarrow{1} A \to 0 \to \ldots
$$
where $A$ is placed in degree $j$ and $j + 1$. Since $\eta_f$ is
compatible with direct sums, we see that the map
$$
(1, d^i) : (\eta_fN)^i / f(\eta_fN)^i \longrightarrow
f^iN^i/f^{i + 1}N^i \oplus f^{i + 1}N^{i + 1}/f^{i + 2}N^{i + 1}
$$
is the direct sum of the corresponding map for
$M^\bullet$ and for $Q^\bullet$. By the universal property
defining the ideals in question, we conclude that
$J_i(N^\bullet, f) = J_i(M^\bullet, f) + J_i(Q^\bullet, f)$.
Hence it suffices to show that $J_i(Q^\bullet, f) = 0$
for all $i$. This is a computation that we omit.
\end{proof}

\end{document}
